\section{Gamma-weighted harmonic identities and differential ladders}
\label{sec:gamma}

\subsection{The complete-harmonic weight at every integer order}
For $q\ge1$ define the local harmonic increment
\begin{equation}\label{eq:local-harmonic}
 h_q(n,a)=\sum_{j=0}^d\frac1{(n+a+j)^q}.
\end{equation}
The complete exponential Bell polynomials are normalized by
\begin{equation}\label{eq:bell}
 \exp\left(\sum_{j\ge1}x_j\frac{u^j}{j!}\right)
 =\sum_{k\ge0}\Bell_k(x_1,\ldots,x_k)\frac{u^k}{k!},
 \qquad \Bell_0=1.
\end{equation}

\begin{theorem}[All-particle Gamma--harmonic evaluation]\label{thm:gamma}
For every integer $m\ge1$,
\begin{equation}\label{eq:harmonic-weight}
 w_n(m,a)=\frac{\Gamma(n+a)}{\Gamma(n+a+d+1)}
 \frac{\Bell_{m-1}(h_1,1!h_2,\ldots,(m-2)!h_{m-1})}{(m-1)!}.
\end{equation}
For $m=1$ the Bell factor is interpreted as one. For $m\ge2$, the
ordered $r$-particle sum of the products of these weights satisfies
\begin{align}
 &\sum_{0\le n_1<\cdots<n_r}\Delta(\bm n)^2
 \prod_{i=1}^r
 \left\{\frac{\Gamma(n_i+a)}{\Gamma(n_i+a+d+1)}
 \frac{\Bell_{m-1}(h_1(n_i,a),\ldots,(m-2)!h_{m-1}(n_i,a))}
 {(m-1)!}\right\}\notag\\
 &\hspace{22mm}=U_r(m,a)+
        \frac{(-1)^m}{(m-1)!}V_r(m,a)\psi^{(m-1)}(a).
 \label{eq:gamma-evaluation}
\end{align}
At positive integral $a$, all the $h_q$ are differences of ordinary
generalized harmonic numbers. At rational $a$, they are rational-grid
harmonic increments.
\end{theorem}
\begin{proof}
The elementary beta partial-fraction identity gives
\[
 w_n(1,a)=\frac1{(n+a)_{d+1}}
        =\frac{\Gamma(n+a)}{\Gamma(n+a+d+1)}.
\]
One can also obtain it by integrating
$t^{n+a-1}(1-t)^d/d!$ on $(0,1)$.
Differentiate \eqref{eq:weight} $m-1$ times in $a$:
\begin{equation}\label{eq:weight-derivative-integer}
 w_n(m,a)=\frac{(-1)^{m-1}}{(m-1)!}
          \partial_a^{m-1}w_n(1,a).
\end{equation}
The logarithmic derivatives of the reciprocal rising factorial are
\[
 \partial_a^q\log w_n(1,a)
     =(-1)^q(q-1)!h_q(n,a).
\]
The Bell formula for an exponential derivative, together with its weighted
homogeneity, proves \eqref{eq:harmonic-weight}. Theorem~\ref{thm:single}
and the classical relation
$\zeta(m,a)=(-1)^m\psi^{(m-1)}(a)/(m-1)!$ prove
\eqref{eq:gamma-evaluation} \cite{dlmf:hurwitz}.
\end{proof}

The first weights are particularly transparent:
\begin{align*}
 w_n(2,a)&=\frac{h_1(n,a)}{(n+a)_{d+1}},\\
 w_n(3,a)&=\frac{h_1(n,a)^2+h_2(n,a)}{2(n+a)_{d+1}},\\
 w_n(4,a)&=\frac{h_1(n,a)^3+3h_1(n,a)h_2(n,a)+2h_3(n,a)}
                  {6(n+a)_{d+1}}.
\end{align*}
Thus the theorem evaluates nonlinear harmonic expressions, not only
partial-fraction weights written in an unfamiliar notation.

For $r=2$, set
\[
 L_n=\frac1{n+1}+\frac1{n+2}+\frac1{n+3},\qquad
 P_n=(n+1)(n+2)(n+3).
\]
A concrete identity is
\begin{equation}\label{eq:concrete-harmonic}
 \boxed{\displaystyle
 \sum_{0\le n<k}\frac{(k-n)^2L_nL_k}{P_nP_k}
       =\frac{\pi^2-7}{16}.}
\end{equation}
Replacing each $L_n$ by
$\bigl(L_n^2+(n+1)^{-2}+(n+2)^{-2}+(n+3)^{-2}\bigr)/2$
gives $(14\zeta(3)-15)/32$. At three particles, the local interval
has five terms and the same construction gives
\eqref{eq:r3m2}--\eqref{eq:r3m3}.

\subsection{Spatial derivatives as exact insertion identities}
For complex $s$ and every integer $q\ge0$,
\begin{equation}\label{eq:shift-weight}
 \partial_a^q w_n(s,a)=(-1)^q(s)_q w_n(s+q,a).
\end{equation}
Consequently, on the original convergence domain,
\begin{align}
 \partial_a^q\Z_r(s,a)
 ={}&(-1)^q\sum_{k_1+\cdots+k_r=q}
    \frac{q!}{k_1!\cdots k_r!}\prod_{i=1}^r(s)_{k_i}\notag\\
 &\qquad\times\sum_{n_1<\cdots<n_r}\Delta(\bm n)^2
       \prod_i w_{n_i}(s+k_i,a).
 \label{eq:insertion-ladder}
\end{align}
All the coefficients and sums are explicit. The right side of
\eqref{eq:single} differentiates to
\begin{equation}\label{eq:spatial-zeta}
 \partial_a^q\Z_r(s,a)=\partial_a^q U_r(s,a)
 +\sum_{j=0}^q\binom qj\partial_a^{q-j}V_r(s,a)
          (-1)^j(s)_j\zeta(s+j,a).
\end{equation}
This uses the standard Hurwitz parameter-derivative law, so again only
single zeta coordinates occur.

For example, let
\[
 \mathcal I_r(s,a)=\sum_{n_1<\cdots<n_r}\Delta(\bm n)^2
        \sum_{i=1}^r w_{n_i}(s+1,a)\prod_{j\ne i}w_{n_j}(s,a).
\]
Then
\begin{equation}\label{eq:first-insertion}
 \mathcal I_r(s,a)=V_r(s,a)\zeta(s+1,a)
   -\frac{\partial_aU_r(s,a)+\partial_aV_r(s,a)\zeta(s,a)}s.
\end{equation}
For $\Rea s>1$, an anchored primitive is therefore
\begin{equation}\label{eq:insertion-primitive}
 \int_{a_0}^{a}\mathcal I_r(s,x)\dd x
       =-\frac{\Z_r(s,a)-\Z_r(s,a_0)}s,
 \qquad \Rea a_0,\Rea a>0,
\end{equation}
with the path lying in the right half-plane. The integral is path
independent because the integrand is holomorphic there. Higher insertion
combinations in \eqref{eq:insertion-ladder} have the analogous anchored
primitives obtained by reducing one derivative order. Thus the harmonic
identities come with differential and antiderivative relations, not just
isolated evaluations.

\subsection{Rational shifts and arithmetic coordinates}
For integer $m$, $U_r(m,a)$ and $V_r(m,a)$ are rational functions of $a$.
At $a\in\Q_{>0}$, the all-particle value therefore lies in the
$\Q$-linear span of $1$ and $\zeta(m,a)$.
For instance, the half-shift identity
$\zeta(s,\tfrac12)=(2^s-1)\zeta(s)$ gives ordinary-zeta reductions
of \eqref{eq:single} at $a=1/2$. General rational shifts may instead
be written with residue-class Dirichlet series or finite root filters.
No claim of arithmetic independence of those coordinates is implied.

For noninteger $s$, the finite exponential factors in $U_r,V_r$ remain
well defined. Their continuation should not be confused with a statement
that the corresponding \emph{cutoff generating function} is always a
finite combination of classical polylogarithms. The latter assertion is
proved below only for rational weights, in particular at integer spectral
orders.
